\documentclass[a4paper]{article}
\usepackage{graphicx}
\usepackage{amsmath,amssymb}
\usepackage{hyperref}

\title{Event Specific Knowledge}
\begin{document}
    \maketitle
    
    Event-specific knowledge (ESK) is vividly detailed information about individual events, often in the form of visual images and sensory-perceptual features.[4] The high levels of detail in ESK fade very quickly, though certain memories for specific events tend to endure longer.[9] Originating events (events that mark the beginning of a path towards long-term goals), turning points (events that re-direct plans from original goals), anchoring events (events that affirm an individual's beliefs and goals) and analogous events (past events that direct behaviour in the present) are all event specific memories that will resist memory decay.[9]
    
    The sensory-perceptual details held in ESK, though short-lived, are a key component in distinguishing memory for experienced events from imagined events.[10] In the majority of cases, it is found that the more ESK a memory contains, the more likely the recalled event has actually been experienced.[10] Unlike lifetime periods and general events, ESK are not organized in their grouping or recall. Instead, they tend to simply 'pop' into the mind.[4] ESK is also thought to be a summary of the content of episodic memories, which are contained in a separate memory system from the autobiographical knowledge base.[6] This way of thinking could explain the rapid loss of event-specific detail, as the links between episodic memory and the autobiographical knowledge base are likewise quickly lost.[6]
    \url{https://en.wikipedia.org/wiki/Autobiographical_memory?utm_source=chatgpt.com#Autobiographical_knowledge_base}
\end{document}